% quiz-macros.tex -- shared package setup for quiz-template.tex (exam class).
%
% NOTE ON PATHS: every template in this repository is stored one level up
% from where it is compiled (see README.md). All relative \input paths
% below are resolved against the folder the document is compiled in,
% which sits one directory above this file.

\newcommand{\coursename}{[Course, set with CourseName]}
\newcommand{\CourseName}[1]{\renewcommand{\coursename}{#1}}

\newcommand{\timelimit}{[Time Limit, set with TimeLimit]}
\newcommand{\TimeLimit}[1]{\renewcommand{\timelimit}{#1}}

\newcommand{\semester}{[Semester, set with Semester]}
\newcommand{\Semester}[1]{\renewcommand{\semester}{#1}}

\newcommand{\examtype}{[Quiz or Exam, set with IsQuiz or IsExam]}
\newcommand{\IsQuiz}{\renewcommand{\examtype}{quiz}}
\newcommand{\IsExam}{\renewcommand{\examtype{exam}}}

% bm: bold vectors (\bm{v}) for physics-style notation in quiz math.
\usepackage{bm}

% --- exam-class patch: point labels next to each \item ---------------------------
% Redefines exam.cls's internal \item@points@pageinfo so the "N pts" label
% is typeset in a group and inserted into the \@labels box -- the box
% exam.cls uses (since v2.218) to attach page-info to the right \item even
% when a question starts with a list/verbatim environment (those
% environments clobber \everypar, which is why the old approach broke).
% \makeatletter is required so the @-names below resolve; the matching
% \makeatother below restores normal catcodes for the rest of the document.
\usepackage{graphicx}
\usepackage{setspace}
\usepackage{enumitem}
\usepackage{amsmath, amsfonts}
\usepackage[none]{hyphenat} %prevent hyphenation
\makeatletter
\ExplSyntaxOn
\def\item@points@pageinfo{%
  \item 
  \if@placepoints
    \if@bonus
      \def\padded@point@block{%
        \begingroup
          \@placepointstrue
          \bonuspoint@block
        \endgroup
      }%
    \else
      \def\padded@point@block{%
        \begingroup
          \@placepointstrue
          \point@block
        \endgroup
      }%
    \fi
    \if@pointsdropped
      % Do nothing!
    \else
      \if@bonus
        \if@bonusqformat
          \ifx\ques@ref\@queslevel
            % Do nothing
          \else
            \setup@point@toks
          \fi
        \else
          \setup@point@toks
        \fi
      \else
        \if@qformat
          \ifx\ques@ref\@queslevel
            % Do nothing
          \else
            \setup@point@toks
          \fi
        \else
          \setup@point@toks
        \fi
      \fi
    \fi
    \global \@placepointsfalse
  \fi
  % Version 2.218$\beta$, 2007-10-31 changes:
  % Instead of appending
  %   \the \pageinfo@commands \the \point@toks
  % to \everypar, we insert them into the box \@labels.  This corrects
  % the problem that arose when a question (or part, etc.) begins with
  % a list environment (including verbatim, flushleft, center,
  % flushright, and possibly others that are implemented as trivlist
  % environments).  The \item command in those environments throws
  % away the previous contents of \everypar, and so the tokens \the
  % \pageinfo@commands \the \point@toks didn't get inserted where we
  % expected.  List environments *do* preserve the contents of the box
  % \@labels, though.
  \global\setbox\g__block_labels_box\hbox{\unhbox\g__block_labels_box %<----------use new name ...
    \the \pageinfo@commands
    \the \point@toks}%
}
\ExplSyntaxOff
\makeatother

% --- Tagged PDF: table headers -----------------------------------------------------
% Tell tagpdf to mark row 1 of every table as a header row, so screen
% readers announce it as a table header.
\tagpdfsetup{table/header-rows={1}}

% --- Name / Student ID box --------------------------------------------------------------
% Prints the answer instructions plus fill-in lines for Name and Student
% ID at the top of the quiz sheet.
\usepackage{multicol}
\usepackage{tabularx}
\usepackage{lastpage}
\newcolumntype{L}{>{\raggedright\arraybackslash}X}
\makeatletter
\def\@maketitle{%
  \newpage
  \null
  {%
    \setlength{\parindent}{0pt}%

    \begin{tabularx}{\textwidth}{l X}
      \textbf{\coursename} & \textbf{Name:} \hrulefill \\

      \textbf{\semester} & \\

      \textbf{\@title} & \\

      \textbf{\@date} & \\

      \textbf{Time Limit: \timelimit}  & \textbf{Student ID:} \hrulefill
    \end{tabularx}
  }

  \hrule
  \vfill
  \begin{singlespace}
{
  \setlength{\parindent}{0in}

  This \examtype\ consists of \numquestions\ questions across \pageref{LastPage} pages worth \numpoints\ points.

  \textbf{Rules:}
}
\begin{itemize}[nosep]
  \item You have the indicated amount of time above to complete the \examtype\, and you must not begin working on the \examtype\ until instructed to do so.

  \item You must show all of your work to receive full credit, unless stated otherwise. \textbf{On problems that require work to be shown, answers that are mysterious or unsupported will not receive full credit.}

  \item You may not use a calculator, textbook, notes or any other external resources.

  \item Please write your initials at the top of every page in the event the pages become separated.

  \item If you run out of space on a page, feel free to use the back of the page to continue your work.

  \item You understand that if you violate any of these or the University's rules, you will receive a score of zero on this \examtype\, and you may be subject to further disciplinary action.
\end{itemize}

{
  \setlength{\parindent}{0in}
  \vfill
  Following the University of Minnesota’s policy on academic honesty, I pledge that the work shown on this \examtype\ is entirely my own and that I have followed the rules stated above. 

  \vfill
  Signature: \hrulefill
}
\end{singlespace}

\begin{center}
  \pointtable[v][questions]
\end{center}
}
\makeatother
% 4. Reduce space below the Date (before main text)
% --- Points summary ----------------------------------------------------------------------
% "This quiz is worth \numpoints points ... 15 minutes" plus the exam-
% class \pointtable listing the points per question.
\newcommand{\PointsSection}{
  This quiz is worth \numpoints\ points and consists of \numquestions\ questions. You have 5 minutes to complete the quiz.

  \begin{center}
    \pointtable[v][questions]
  \end{center}
}

% --- Math -------------------------------------------------------------------------------
\usepackage{amsmath,amssymb}

% --- Shared course macros ---------------------------------------------------------------
% Root macros.tex: math shorthands used across all course documents
% (\RR, \NN, \ZZ, \QQ, \grad, \curl, \tr, \der, \pder, ...).
\input{../../macros.tex}
